\subsubsection{Criterios de selección de la Fuente de alimentación}

Para realizar la selección de la potencia de la fuente de alimentación, es necesario repasar cómo se alimentan cada uno de los módulos que integran el sistema para estimar la potencia necesaria para funcionamiento del circuito.

Haciendo un repaso por los diferentes subsistemas del compostador, se puede observar que está compuesto por diversos elementos electrónicos de control, los cuales deben de estar conectados a una tensión de alimentación de 5 V DC a 1 A.

Estos sistemas están constituidos por la tarjeta de desarrollo, diversos circuitos integrados, elementos resistivos e inductivos que se especifican en la comunicación de los diversos subsistemas internos, así como de señales de control para la activación de los diferentes procesos.

Para el sistema de riego, se tiene la activación de una pequeña bomba de agua sumergible, la cual se alimenta a 127 V AC y consume una potencia de 40 W.

También  el sistema de regulación de temperatura cuenta con una resistencia eléctrica que eleva la temperatura. Este elemento funciona con una potencia de 800 W a 127 V AC. 

Este tipo de elementos eléctricos están diseñados para ser conectados directamente a la toma de corriente, por lo que la fuente seleccionada debe considerar una salida que se adapte a estas necesidades o que permita realizar un arreglo para su conexión.

Para el sistema de ventilación, es necesario contar con una salidas para conectar el ventilador que s se encarga de generar el flujo de aire necesario para oxigenar y para bajar la temperatura al interior del contenedor. 

Este elemento está alimentado por una tensión de 12 V DC con una corriente de 300 mA, de esta forma, una salida de corriente directa de 12 V DC a 1 A es suficiente para alimentar todo el sistema.
 
Para los motores de aireacion y el de trituracion, se han considerados unas entradas de corriente de 127 V AC a 5.6 A. que son sufientes para arrancar estos elementos.
